% ============================================================
% EXPERIMENTAL SETUP
% ============================================================
\section{Experimental Setup}

\subsection{Dataset Preparation}

There is no public corpus of faulty medical-data programs, and real clinical scripts could not be shared anyway, so we
generated one. A generator assembles valid programs from three to seven of sixteen building blocks (capped dosing, BMI,
vital-sign alerts, shock index, a sepsis rule, infusion rate, an optional lab value with three guarding styles, PHI
handling, a scoring loop, unit conversion and so on). Two blocks are hard negatives that look wrong but are correct: a
division by a variable whose range includes zero but is guarded by a positive test, and a threshold equal to the upper
bound of the declared range. Names, thresholds, ranges, comments and block order are randomised, and every generated
program must compile cleanly.

Forty-six mutation operators, three to seven per family, then edit a valid program: deleting a semicolon (syntax),
swapping \code{degC} for \code{mmHg} (unit), printing a PHI input directly (phi), removing a \code{present} guard
(missing), inverting an alert comparison (logic), and so on. A mutated program is kept only if the compiler reports a
diagnostic of the intended family. The exception is the two ``soft'' logic operators (inverting a rule, or giving it the
alert text of the opposite kind), which by design produce no diagnostic. In about 20\% of cases a second, lower-priority
fault is injected first, so the program has two defects but one well-defined label.

The corpus has \NProg{} unique programs, 1{,}000 per class, averaging \MeanLoc{} lines. Of the 9{,}000 faulty ones,
\NDouble{} (\DoubleShare\%) contain two faults. A fixed stratified split holds out 200 programs per class as the
\emph{test set} (\NTest{} programs), which is never used for training, GAN fitting or tuning. The other 800 per class form
the training pool. Generating the whole corpus, including the compile check on each fault, took 18 seconds.

All experiments ran on a laptop CPU (Intel Core i5-13420H, 16 GB RAM, Windows 11) with Python 3.13, scikit-learn 1.9,
PyTorch 2.14 and NumPy 2.5. The networks are small and the vectors have 101 dimensions, so no GPU was needed.

\subsection{Experimental Protocol}

All random choices are seeded.

\begin{itemize}
    \item \textbf{E1, rules versus learning.} The four classifiers are trained on all 8{,}000 pool programs and tested on
    the 2{,}000 test programs, next to the two rule baselines.
    \item \textbf{E2, feature views.} A random forest is trained on four feature subsets with 5 to 800 programs per class,
    five random draws each.
    \item \textbf{E3, augmentation under scarcity.} The valid class has 400 real programs and each faulty class only
    $n\in\{5,10,20,50,100\}$. Seven regimes are compared: no augmentation, class weights, random oversampling, Gaussian
    jitter (noise at 10\% of the feature s.d.), SMOTE with $k=5$, and the BCE and WGAN-GP GANs. All but the first two bring
    each faulty class to 400 vectors. Every setting runs over ten seeds and is tested with a random forest and an MLP.
    Differences are checked with a paired $t$-test across seeds.
    \item \textbf{E4, anomaly detection.} Detectors are trained on the 800 valid pool programs and tested on 200 valid
    versus 200 faulty test programs per family.
    \item \textbf{E5, compiler performance.} Valid programs of growing size are compiled five times each (median time per
    phase), and 300 programs are checked for behavioural equivalence before and after optimisation on ten random inputs.
\end{itemize}

Classifiers are scored with accuracy and macro-F1 \cite{sokolova2009}, where macro-F1 is the mean over classes of
$F1_k=2P_kR_k/(P_k+R_k)$. On a balanced test set the two are close. Detectors are scored with AUROC, the probability that a
random faulty program scores higher than a random valid one~\cite{fawcett2006}.
